\documentclass[12pt]{article}
\usepackage[T1]{fontenc}
\usepackage[utf8]{inputenc}
\usepackage{lmodern}
\usepackage{textcomp}
\usepackage{enumitem}
\usepackage{geometry}
\geometry{margin=1in}
\begin{document}
\begin{center}
Answer Key - Chemistry - Undergraduate Level Assessment 6 \
\small (Answers are ordered exactly as in the question paper.)
\end{center}

\section*{Multiple Choice}
\begin{enumerate}[label=\arabic*.]
\item \textbf{Answer:} D
\\ \textit{Options:} A) Ammonium hydroxide; B) Sulfuric acid; C) Acetic acid; D) Potassium hydrogen phthalate
\\ \textit{Note:} Potassium hydrogen phthalate is a primary standard because it is a stable, pure compound with a known purity that can be accurately weighed and dissolved to create a solution of known concentration for standardizing bases in titrations.
\item \textbf{Answer:} A
\\ \textit{Options:} A) Blackbody radiation curves; B) Emission spectra of diatomic molecules; C) Electron diffraction patterns; D) Temperature dependence of reaction rates
\\ \textit{Note:} Planck's quantum theory successfully explains blackbody radiation curves by introducing the concept of quantized energy levels, which resolved the ultraviolet catastrophe and accurately described the observed spectral distribution of electromagnetic radiation emitted by idealized black bodies.
\item \textbf{Answer:} D
\\ \textit{Options:} A) BaH 2; B) BaOH; C) Ba; D) BaO
\\ \textit{Note:} The anhydride of Ba(OH)2 is BaO because it represents the product formed when barium hydroxide loses water, resulting in the formation of barium oxide, which lacks hydroxide ions.
\item \textbf{Answer:} B
\\ \textit{Options:} A) 0.0471 Hz; B) 21.2 Hz; C) 42.4 Hz; D) 66.7 Hz
\\ \textit{Note:} The linewidth can be calculated using the formula Δν = 1/(2πT 2), and substituting T$_{2}$ = 15 ms yields a linewidth of approximately 21.2 Hz.
\item \textbf{Answer:} B
\\ \textit{Options:} A) Ti$_{3}$+; B) Cr$_{3}$+; C) Fe$_{3}$+; D) Zn$_{2}$+
\\ \textit{Note:} Cr$_{3}$+ is a d 3 metal ion that has all of its unpaired electrons paired, resulting in a lack of net magnetic moment, which disqualifies it from acting as a paramagnetic quencher.
\end{enumerate}

\section*{True / False}
\begin{enumerate}[label=\arabic*.]
\setcounter{enumi}{5}
\item \textbf{Answer:} True
\\ \textit{Options:} A) True; B) False
\\ \textit{Note:} H$_{3}$O+ is the conjugate acid of H$_{2}$O, which acts as its conjugate base, demonstrating the relationship between acids and bases in the Brønsted-Lowry acid-base theory.
\item \textbf{Answer:} False
\\ \textit{Options:} A) True; B) False
\\ \textit{Note:} BeCl 2 is more likely to act as a Lewis acid compared to ZnCl 2 because beryllium has an empty p orbital that can accept electron pairs, while zinc, despite being in a higher oxidation state, has a filled d orbital that does not favor Lewis acid behavior.
\item \textbf{Answer:} True
\\ \textit{Options:} A) True; B) False
\\ \textit{Note:} The statement is true because the principal allotropes of elemental boron, such as α-rhombohedral boron and β-boron, feature B$_{12}$ icosahedra as their fundamental structural units, which contribute to their unique properties and stability.
\item \textbf{Answer:} False
\\ \textit{Options:} A) True; B) False
\\ \textit{Note:} The statement is false because the 13 C spectrum of hexane shows five distinct chemical shifts due to its symmetrical structure resulting in equivalent carbon environments, not because it has multiple isomers.
\item \textbf{Answer:} False
\\ \textit{Options:} A) True; B) False
\\ \textit{Note:} The limiting high-temperature molar heat capacity at constant volume (C\_V) of a gas-phase diatomic molecule is actually 5/2 R, not 2.5 R, because it accounts for translational and rotational degrees of freedom.
\end{enumerate}

\section*{Long Form Response}
\begin{enumerate}[label=\arabic*.]
\setcounter{enumi}{10}
\item \textbf{Answer:} The atomic radii of lanthanide elements exhibit a general trend of decreasing size from lanthanum (La) to lutetium (Lu). This trend can be attributed to the increasing effective nuclear charge experienced by the outer electrons as more protons are added to the nucleus while the additional f-electrons do not shield the nuclear charge effectively. Consequently, the increased attraction between the nucleus and the outermost electrons results in a smaller atomic radius. A misleading statement regarding this trend may suggest that atomic radii increase across the series, which contradicts the observed decrease due to the aforementioned reasons. Thus, the correct understanding of the trend is crucial for comprehending the properties of these elements.
\\ \textit{Note:} correct because it accurately describes the trend of decreasing atomic radii in lanthanides due to the increasing effective nuclear charge from added protons, while the f-electrons provide inadequate shielding, leading to a stronger attraction between the nucleus and outer electrons.
\item \textbf{Answer:} In thermodynamics, spontaneity refers to a process that occurs without external intervention, and its relationship with entropy is fundamental. According to the second law of thermodynamics, for a process to be spontaneous, the total entropy of the system and its surroundings must increase. This increase in total entropy signifies that energy disperses and systems evolve toward a state of greater disorder. Essentially, spontaneous processes tend to favor configurations that maximize the number of accessible microstates, thus enhancing the overall entropy. Therefore, the requirement for total entropy increase is a criterion that ensures the natural progression of thermodynamic processes toward equilibrium.
\\ \textit{Note:} correct because it accurately describes the principle that spontaneous processes are characterized by an increase in total entropy, reflecting the natural tendency of systems to evolve towards greater disorder and energy dispersion in accordance with the second law of thermodynamics.
\end{enumerate}

\vfill
\noindent\textit{End of Answer Key}
\end{document}